\documentclass[11pt,a4paper]{article}

% ---------- packages ----------
\usepackage[utf8]{inputenc}
\usepackage[T1]{fontenc}
\usepackage{lmodern}
\usepackage[margin=2.2cm]{geometry}
\usepackage{xcolor}
\usepackage[most]{tcolorbox}
\usepackage{titlesec}
\usepackage{enumitem}
\usepackage{booktabs}
\usepackage{longtable}
\usepackage{array}
\usepackage{fancyhdr}
\usepackage{listings}
\usepackage[hidelinks]{hyperref}

% ---------- palette ----------
\definecolor{spInk}{HTML}{1F2A44}      % near-black ink
\definecolor{spPrimary}{HTML}{2A4D9B}  % indigo
\definecolor{spAccent}{HTML}{0E7C7B}   % teal
\definecolor{spHow}{HTML}{1565C0}      % blue
\definecolor{spWhy}{HTML}{2E7D32}      % green
\definecolor{spSec}{HTML}{B23A00}      % burnt orange (security)
\definecolor{spQA}{HTML}{5E35B1}       % purple (Q&A)
\definecolor{spGrayBg}{HTML}{F4F6FB}
\definecolor{spRule}{HTML}{C9D2E8}

\hypersetup{colorlinks=true,linkcolor=spPrimary,urlcolor=spAccent}

% ---------- section styling ----------
\titleformat{\section}{\Large\bfseries\color{spPrimary}}{\thesection.}{0.5em}{}
\titleformat{\subsection}{\large\bfseries\color{spAccent}}{\thesubsection}{0.5em}{}
\titlespacing*{\section}{0pt}{14pt}{6pt}

% ---------- header / footer ----------
\pagestyle{fancy}\fancyhf{}
\renewcommand{\headrulewidth}{0.4pt}
\fancyhead[L]{\small\color{spPrimary}\textbf{ScholarPath}}
\fancyhead[R]{\small\color{spInk}Technical Architecture \& Security Brief}
\fancyfoot[C]{\small\thepage}

% ---------- colored callout boxes ----------
\newtcolorbox{howbox}{breakable,colback=spHow!4,colframe=spHow!75!black,
  boxrule=0.7pt,arc=2pt,left=6pt,right=6pt,top=4pt,bottom=4pt,
  fonttitle=\bfseries\color{white},colbacktitle=spHow!85!black,title=How it works}
\newtcolorbox{whybox}{breakable,colback=spWhy!4,colframe=spWhy!70!black,
  boxrule=0.7pt,arc=2pt,left=6pt,right=6pt,top=4pt,bottom=4pt,
  fonttitle=\bfseries\color{white},colbacktitle=spWhy!75!black,title=Why we did it this way}
\newtcolorbox{secbox}{breakable,colback=spSec!5,colframe=spSec!75!black,
  boxrule=0.7pt,arc=2pt,left=6pt,right=6pt,top=4pt,bottom=4pt,
  fonttitle=\bfseries\color{white},colbacktitle=spSec!80!black,title=Security rationale}
\newtcolorbox{qabox}{breakable,colback=spQA!4,colframe=spQA!70!black,
  boxrule=0.7pt,arc=2pt,left=6pt,right=6pt,top=4pt,bottom=4pt,
  fonttitle=\bfseries\color{white},colbacktitle=spQA!75!black,title=Likely question}

% ---------- code listing style ----------
\lstset{basicstyle=\ttfamily\small,backgroundcolor=\color{spGrayBg},
  frame=single,rulecolor=\color{spRule},breaklines=true,
  keywordstyle=\color{spPrimary}\bfseries,commentstyle=\color{spAccent}\itshape,
  columns=fullflexible,showstringspaces=false,xleftmargin=4pt,framexleftmargin=2pt}

\newcommand{\code}[1]{\texttt{\textcolor{spInk}{#1}}}
\setlist{nosep,leftmargin=1.4em,topsep=2pt}

% ====================================================================
\begin{document}

% ---------- title block ----------
\begin{tcolorbox}[colback=spPrimary,colframe=spPrimary,arc=4pt,boxrule=0pt]
  \color{white}
  {\Huge\bfseries ScholarPath}\\[2pt]
  {\LARGE Technical Architecture \& Security Brief}\\[6pt]
  {\normalsize A walkthrough of the architecture, the request pipeline (CQRS),
  data integrity, and the security-critical features --- encryption, file
  uploads \& scanning, PII handling, live meetings, AI/RAG, and the
  booking--payment money flow. Every mechanism below is taken from the live
  codebase, not a plan.}
\end{tcolorbox}

\vspace{4pt}
{\small\color{spInk}\textbf{Scope:} .NET 10 Clean-Architecture monolith $\cdot$
React 19 SPA $\cdot$ Azure $\cdot$ $\approx$55 entities $\cdot$ $\approx$33 ports /
$\approx$46 adapters $\cdot$ bilingual EN/AR.}

\vspace{2pt}\hrule height 0.4pt\vspace{6pt}
{\small\tableofcontents}
\vspace{4pt}\hrule height 0.4pt

% ====================================================================
\section{Architecture --- the foundation}

ScholarPath is a \textbf{Clean-Architecture} monolith. Dependencies always point
\emph{inward}; the inner layers know nothing about the outer ones.

\begin{center}\small\ttfamily
API \ $\rightarrow$ \ Application \ $\rightarrow$ \ Domain \qquad
Infrastructure \ $\rightarrow$ \ (implements Application \& Domain ports)
\end{center}

\begin{itemize}
\item \textbf{Domain} --- entities, enums, domain events, and two domain ports
  (\code{ICurrentUserService}, \code{IDateTimeService}). No framework dependency.
\item \textbf{Application} --- use-cases as \textbf{CQRS} commands/queries (MediatR),
  FluentValidation validators, AutoMapper profiles, the AI cost gate, and the
  $\approx$33 \textbf{ports} (interfaces) the outside world must satisfy.
\item \textbf{Infrastructure} --- the $\approx$46 \textbf{adapters} implementing
  those ports: EF Core, Stripe, Azure OpenAI, ACS, Blob, Key Vault, Event Hub,
  SignalR, Hangfire jobs.
\item \textbf{API} --- controllers + Stripe webhook receiver, ASP.NET Identity +
  JWT, three SignalR hubs.
\end{itemize}

\begin{whybox}
The \emph{ports \& adapters} seam is the most important decision in the system.
It makes the code (a)~\textbf{testable} --- swap a real service for a stub;
(b)~\textbf{portable} --- swap a provider by \emph{configuration}, with no code
change (\code{Ai:Provider}, \code{Storage:Provider}, \code{Email:Provider},
\code{Acs:ConnectionString}, \code{FileScanning:Enabled}, \code{Hangfire:Enabled});
and (c)~keeps business rules independent of Azure.
\end{whybox}

% ====================================================================
\section{CQRS \& the request pipeline}

Every use-case is a \textbf{command} (writes/mutations) or a \textbf{query}
(reads), dispatched through \textbf{MediatR}. This is CQRS: the read path and the
write path are separate request types with separate handlers and validators.

\subsection{The behavior pipeline (order matters)}
Every request passes through four \emph{pipeline behaviors} before reaching its
handler. They are registered --- and therefore execute --- in this exact order
(outermost first):

\begin{lstlisting}
cfg.AddOpenBehavior(typeof(LoggingBehavior<,>));      // 1 outermost
cfg.AddOpenBehavior(typeof(ValidationBehavior<,>));   // 2
cfg.AddOpenBehavior(typeof(PerformanceBehavior<,>));  // 3
cfg.AddOpenBehavior(typeof(AuditBehavior<,>));        // 4 closest to handler
\end{lstlisting}

\begin{enumerate}
\item \textbf{Logging} (outermost) --- wraps everything, so it records the request
  even if validation rejects it.
\item \textbf{Validation} --- runs all FluentValidation validators for the request
  and short-circuits with a 400 before any work happens. Placed early on purpose:
  no point timing or auditing a request that is structurally invalid.
\item \textbf{Performance} --- times the handler and flags slow requests.
\item \textbf{Audit} (innermost) --- records the actual business action right
  around the handler, where the request is known-valid.
\end{enumerate}

\begin{whybox}
Ordering is deliberate: \textbf{Logging} must see every request; \textbf{Validation}
must reject bad input \emph{before} we spend time or write audit rows;
\textbf{Audit} must sit closest to the handler so it only records actions that
actually executed. Validators, AutoMapper profiles, and the refund calculator are
auto-discovered from the assembly, so adding a feature needs no DI wiring.
\end{whybox}

% ====================================================================
\section{Data integrity \& ``non-appearance'' (soft delete)}

\subsection{Soft delete --- how a row ``disappears'' without being destroyed}
Sensitive/business rows are never hard-deleted. They implement
\code{ISoftDeletable} (\code{IsDeleted, DeletedAt, DeletedByUserId}); ``deleting''
flips \code{IsDeleted = true}. To make those rows vanish from the whole
application automatically, each such entity has an EF Core \textbf{global query
filter}:

\begin{lstlisting}
b.HasQueryFilter(x => !x.IsDeleted);   // applied on ~15 entities
\end{lstlisting}

\begin{howbox}
Once the filter is on an entity, \emph{every} LINQ query EF generates for it gets
an implicit \code{WHERE IsDeleted = 0}. A developer cannot forget to add it ---
soft-deleted rows simply never appear in any query, join, or count, unless code
explicitly opts out with \code{IgnoreQueryFilters()} (used only by admin/GDPR
tooling). 15 entities carry this filter today.
\end{howbox}

\begin{whybox}
Hard deletes would cascade-break history (reviews, payments, audit trails) and
make ``who deleted what, when'' impossible. Soft delete keeps referential history
intact, supports an undo/restore path, and --- combined with the audit log ---
gives full traceability. It is also the safe default for a platform that must
honor data-retention and dispute rules.
\end{whybox}

\subsection{The \code{SaveChangesAsync} interceptor (one place, many guarantees)}
The \code{DbContext} overrides \code{SaveChangesAsync} to enforce four
cross-cutting rules in a single choke point:

\begin{itemize}
\item \textbf{Audit stamping} --- sets \code{CreatedAt} on new auditable rows and
  \code{UpdatedAt} on every modification (seeders keep intentional back-dates).
\item \textbf{Domain events after commit} --- domain events are dispatched
  \emph{only after} the transaction is durably persisted, so a side effect never
  fires for a change that rolled back.
\item \textbf{Conflict translation} --- a unique-index violation (SQL 2601/2627) is
  caught and re-thrown as a friendly \code{ConflictException} (e.g.\ ``that slot
  was just booked'') instead of a raw 500.
\item \textbf{Global cascade safety} --- every FK pointing at \code{ApplicationUser}
  is swept to \code{DeleteBehavior.Restrict}, avoiding SQL Server's
  multiple-cascade-path error and protecting user history.
\end{itemize}

\subsection{Other integrity controls}
\begin{itemize}
\item \textbf{Optimistic concurrency} --- \code{RowVersion} on auditable entities
  prevents lost updates.
\item \textbf{Idempotency keys} --- \code{Payment.IdempotencyKey} and
  \code{Notification.IdempotencyKey} make money and messaging exactly-once under
  retries; Stripe webhooks dedupe on \code{StripeWebhookEvent.StripeEventId}.
\item \textbf{Filtered unique indexes} encode real invariants in the database:
  one \emph{active} application per student per scholarship (FR-057), one
  \emph{active} booking per consultant per slot, one \emph{active} financial
  config per type.
\end{itemize}

% ====================================================================
\section{The security model (defense in depth)}

Independent, layered controls so no single failure is catastrophic:

\renewcommand{\arraystretch}{1.25}
\begin{longtable}{>{\bfseries\color{spPrimary}}p{3.0cm}p{11.5cm}}
\toprule
Layer & Control \\ \midrule
Transport & HTTPS everywhere; JWT bearer on the API; authenticated SignalR hubs \\
Identity & ASP.NET Core Identity --- hashed passwords, lockout, email confirmation \\
Tokens & Short-lived access JWT + \emph{rotating} refresh tokens, stored \emph{hashed} \\
Authorization & Role-based (Student/Company/Consultant/Admin) + per-resource ownership \\
Secrets \& keys & Azure \textbf{Key Vault} key providers (field-encryption key, JWT signing key) \\
Data at rest & \textbf{AES-256-GCM} field encryption on sensitive columns \\
Uploads & \textbf{Scan-before-store} with ClamAV, \textbf{fail-closed} \\
Privacy / PII & Prompt redaction before AI; loose references; GDPR export/delete \\
Integrity & Idempotency keys, filtered unique indexes, \code{RowVersion} \\
Auditability & Append-only \code{AuditLog} (before/after JSON); per-call \code{AiInteraction} log \\
Soft delete & Global query filters hide deleted rows from every query \\
\bottomrule
\end{longtable}

% ====================================================================
\section{Encryption at rest (AES-256-GCM)}
Port \code{IFieldEncryptionService} $\rightarrow$ \code{AesGcmFieldEncryptionService},
wired onto chosen PII columns (e.g.\ a profile \code{Biography}, an application's
\code{PersonalNotes}) through an EF Core value converter --- plaintext in C\#,
ciphertext in the database.

\begin{howbox}
AES-256-GCM is an \emph{authenticated} (AEAD) cipher: 256-bit key, a fresh random
12-byte nonce per encryption, and a 16-byte authentication tag. The stored value
is a self-describing envelope
\code{enc:v1:} + Base64(\,nonce $\Vert$ tag $\Vert$ ciphertext\,).
On decrypt the tag is verified --- a flipped bit, a truncation, or the wrong key
makes decryption \emph{throw}, not return garbage. The random nonce means the same
plaintext encrypts to different ciphertext each time.
\end{howbox}

\begin{itemize}
\item \textbf{Versioned} (\code{v1}): a future key rotation introduces \code{v2}
  while still decrypting old \code{v1} rows.
\item \textbf{Safe rollout}: a value without the prefix (legacy plaintext) is
  returned unchanged and re-encrypted on its next write --- zero downtime.
\item \textbf{Key} comes from \code{IFieldEncryptionKeyProvider} (Key Vault in prod,
  local in dev) and must be exactly 32 bytes (checked at start-up).
\end{itemize}

\begin{secbox}
GCM gives confidentiality \emph{and} integrity in one pass; a plain cipher would
protect secrecy but not detect tampering. The key never lives in source or config.
This protects the value end-to-end --- even a DB backup or a DBA sees only
ciphertext for those columns (which transparent DB-level encryption alone cannot
guarantee once a query returns a row).
\end{secbox}

\subsection{Key management}
Two independent providers, each with a prod + dev adapter:
\code{IFieldEncryptionKeyProvider} (KeyVault / Local) for the AES key, and
\code{IJwtKeyProvider} (KeyVault / Local-PEM) for the JWT signing key. Surfacing
the signing key behind its own port is what makes \textbf{key rotation} an
operations action, not a redeploy.

% ====================================================================
\section{Authentication \& authorization}
\begin{itemize}
\item \textbf{ASP.NET Core Identity} --- \code{ApplicationUser : IdentityUser<Guid>};
  passwords hashed by \code{IdentityPasswordHasher}; lockout; email confirmation.
\item \textbf{Tokens} --- \code{TokenService} issues a short-lived access JWT plus a
  refresh token; refresh tokens are \emph{rotated} on use, individually/globally
  revocable, and stored \emph{hashed} (\code{RefreshToken.TokenHash}). Same hashing
  for \code{PasswordResetToken} --- a DB leak yields no usable token.
\item \textbf{SSO} --- \code{SsoService} exchanges OAuth codes with Google and
  Microsoft.
\item \textbf{Authorization} --- role-based (the overlapping Student/Company/
  Consultant/Admin specialization via the \code{UserRoles} join) plus per-resource
  ownership checks in handlers.
\end{itemize}

% ====================================================================
\section{File upload pipeline}
Port \code{IBlobStorageService} $\rightarrow$ \code{FileStorageService}: one
interface, two providers chosen by \code{Storage:Provider} --- \code{Local}
(disk, dev) or \code{AzureBlob} (Azure Storage). Objects use an opaque,
provider-tagged path \code{provider:container/key}, where \code{key} is
GUID-prefixed so same-named uploads never collide; Azure containers are created
\code{PublicAccessType.None} --- never public, downloads always go through the
authenticated API.

\begin{howbox}
\textbf{Scan before store.} The upload handler runs the file through the malware
scanner \emph{while it is still an in-memory stream}; \emph{only a Clean verdict is
persisted}. This guards the document vault (application attachments, upgrade
files, forum attachments, session recordings).
\end{howbox}

% ====================================================================
\section{Malware scanning (ClamAV, fail-closed)}
Port \code{IFileScanService} $\rightarrow$ \code{ClamAvFileScanService}: bytes are
streamed to a ClamAV \code{clamd} daemon over INSTREAM (via \code{nClam}) before
storage. Verdicts: \textbf{Clean} (proceed; stream rewound so the same bytes are
stored), \textbf{Infected} (rejected; signature logged), \textbf{ScanUnavailable}
(rejected).

\begin{secbox}
\textbf{Fail-closed.} Any daemon unreachability, timeout, or error maps to
\code{ScanUnavailable} and the upload is \textbf{rejected} --- an unscanned file is
never stored. (Feature-flagged by \code{FileScanning:Enabled}; \code{NoOpFileScanService}
in dev.) Scanning the stream first, fail-closed, removes the window in which a
malicious file could be reachable.
\end{secbox}

% ====================================================================
\section{PII handling}
PII protection is layered across the AI path, the data model, and the lifecycle.

\begin{howbox}
\textbf{Prompt redaction before persistence \emph{and} before any provider.} In
the chatbot handler the user's message passes through \code{RedactPii()} first:
emails $\rightarrow$ \code{[redacted-email]}, 13--19-digit runs (cards)
$\rightarrow$ \code{[redacted-card]}, 10+-digit runs (phones) $\rightarrow$
\code{[redacted-phone]}. The \emph{redacted} text is what we store
(\code{AiInteraction.PromptText}) and what we send to OpenAI --- the raw message
never leaves the process. (Regex order matters: email, then card, then the most
permissive phone pattern last.)
\end{howbox}

\begin{itemize}
\item \textbf{Personalization without contact PII} --- the student-profile block
  injected for tailored answers carries only academic facts (nationality, level,
  field, GPA, preferred countries), never email/phone.
\item \textbf{Per-user AI memory} --- replayed history is filtered to the requesting
  user, so a leaked \code{SessionId} cannot leak another user's transcript.
\item \textbf{Redaction QA} --- \code{AiRedactionAuditSample} stores periodic samples
  of redacted prompts for human review (a sampling job) --- we can prove the
  redactor works.
\item \textbf{Loose references} --- high-volume / audit / analytics tables (payments,
  notifications, votes, chat, audit log, AI interactions) hold the user
  \code{Guid} with \emph{no FK constraint}, so a user can be deleted or anonymized
  \emph{without} cascade-breaking history.
\item \textbf{GDPR} --- \code{UserDataRequest} drives export and delete
  (\code{DataExportJob}, \code{DataDeleteJob}); sensitive free-text is encrypted at
  rest; profile fields are role-gated in the UI.
\end{itemize}

% ====================================================================
\section{Live meetings (Azure Communication Services)}
Port \code{IMeetingService} $\rightarrow$ \code{AzureCommunicationMeetingService}
(selected when \code{Acs:ConnectionString} is set; \code{StubMeetingService} in dev).

\begin{howbox}
A booking's ``room'' is an \textbf{ACS group-call id (a GUID)} --- no server-side
room object. On join, the API mints a \textbf{short-lived, VoIP-scoped access
token per user} from ACS Identity (\code{CreateUserAndTokenAsync}, with an
expiry), handed only to the two booking participants through the authorized join
endpoint. \textbf{Recording} uses ACS Call Automation
(\code{StartRecordingAsync} $\rightarrow$ \code{RecordingId};
\code{DownloadRecordingAsync} streams it), persisted as a \code{SessionRecording}
+ Blob object.
\end{howbox}

\textbf{Attendance \& no-shows:} the booking records \code{RecordingStartedAt}
(PB-006) and \code{StudentJoinedAt}\,/\,\code{ConsultantJoinedAt} (FR-217); the
\code{MeetingNoShowSweepJob} marks \code{NoShowMarkedAt} and drives the
no-show/refund logic.

\begin{secbox}
A video token is a bearer credential. Minting one \emph{per participant}, scoped
to VoIP and expiring, means a leaked token is useless after the session and cannot
be replayed for another call.
\end{secbox}

% ====================================================================
\section{The AI subsystem}
Port \code{IAiService} has \textbf{three interchangeable adapters} chosen by
\code{Ai:Provider}: \code{AzureOpenAiService}, \code{OpenAiService} (OpenAI direct),
\code{LocalAiService} (deterministic, offline). Embeddings mirror this. Three
operations: \code{Recommend()}, \code{CheckEligibility()}, \code{Ask()}.

\begin{itemize}
\item \textbf{Recommendations are not an LLM call} --- they are produced by the
  deterministic \code{LocalAiService}, scoring scholarships against the student
  profile + metadata (+ RAG), returning the top-N (default 5). The LLM backs
  \emph{chat / Q\&A}.
\item \textbf{Chat (\code{Ask})} carries session memory (last 20 turns,
  user-scoped), an injected profile system-turn for personalization, RAG
  grounding, and returns the answer + a disclaimer + the \emph{sources} used.
\end{itemize}

\begin{secbox}
\textbf{Cost control --- \code{AiCostGate}.} Before \emph{every} AI call we enforce a
rolling 24-hour per-user budget (\code{DailyUserCostLimitUsd}, default \$1.00): we
sum the user's \code{AiInteraction.CostUsd} over 24h and reject the call if this
one would exceed the cap. Every call is logged (provider, model, tokens, cost, and
a metadata blob of RAG sources) --- full spend observability. What keeps the bill
bounded: the per-user gate \emph{and} the fact that the high-frequency feature
(recommendations) is local, not an LLM call.
\end{secbox}

% ====================================================================
\section{RAG --- Retrieval-Augmented Generation}
\begin{howbox}
\textbf{Index} --- \code{KnowledgeBaseIndexer} turns scholarships, resources,
consultants and top community posts into \code{KnowledgeDocument} rows (bilingual
content, a \code{ContentHash}, and an embedding \code{byte[]} tagged with the model
name + dimensions; \code{SourceType+SourceKey} unique, so re-indexing upserts).
\textbf{Retrieve} --- \code{KnowledgeRetriever} embeds the query, ranks the corpus
by \textbf{cosine similarity} and returns the \textbf{top-K}, considering only
documents embedded with the \emph{currently active} model (vectors from another
model live in a different space).
\end{howbox}

\begin{whybox}
The corpus is a few hundred rows, so an in-memory similarity scan per query is
more than fast enough --- we deliberately avoid operating a separate vector
database. Which documents grounded each answer is recorded in
\code{AiInteraction.MetadataJson} and surfaced as ``sources'', giving a RAG audit
trail and reducing hallucination.
\end{whybox}

% ====================================================================
\section{Booking \& the payment money flow}
Paid consultant sessions use an \textbf{authorize-then-capture (escrow-style)}
model on Stripe, so the student's money is \emph{held} but not taken until the
consultant accepts. Port \code{IStripeService} $\rightarrow$ \code{StripeService}
(real) / \code{StubStripeService} (dev).

\subsection{Lifecycle (step by step)}
\begin{enumerate}
\item \textbf{Book \& authorize} --- the student requests a slot; we create a Stripe
  \code{PaymentIntent} with \textbf{manual capture} and \code{PaymentMethodTypes =
  [card]}. The card is \emph{authorized} (funds held), not charged. A \code{Payment}
  row is created with an \code{IdempotencyKey} (the same key is sent to Stripe, so a
  retry never double-charges). Amounts are integer \textbf{cents}.
\item \textbf{Accept $\rightarrow$ capture} --- when the consultant accepts, we
  \textbf{capture} the held intent (\code{CapturePaymentIntentAsync}). A capture on
  an un-authorized intent is surfaced as a clean 422 (``finish payment first''),
  not a 500.
\item \textbf{Reject / expire $\rightarrow$ release} --- the hold is cancelled
  (\code{CancelPaymentIntentAsync}); the student is never charged.
\item \textbf{No-show / cancellation $\rightarrow$ refund} --- a
  \code{RefundCalculatorService} computes the owed amount per policy, then
  \code{RefundPaymentAsync} issues a full/partial refund. \code{MeetingNoShowSweepJob}
  triggers this automatically.
\item \textbf{Payout to consultant} --- consultants onboard a \textbf{Stripe Connect
  Express} account (\code{CreateConnectAccountAsync} + onboarding link). The
  \code{StripePayoutJob} pays out their net via \code{CreatePayoutAsync} (idempotent),
  recording a \code{Payout} row that aggregates the included payment ids.
\end{enumerate}

\subsection{Profit-share split}
Each \code{Payment} stores \code{AmountCents}, \code{ProfitShareAmountCents}
(platform fee) and \code{PayeeAmountCents} (consultant net). The split percentages
come from \code{ProfitShareConfig} / \code{FinancialConfigRule}, each protected by a
\textbf{filtered unique index} that enforces exactly one \emph{active} rule per
payment type (PB-014 / FR-170) --- so fees are configurable and versioned over
time without ambiguity.

\subsection{Webhook reconciliation}
\code{StripeService.ParseWebhook} \textbf{verifies the Stripe signature} and
extracts the typed event (PaymentIntent, Charge, Account, Payout, Dispute). The
handler reconciles the matching \code{Payment} by intent id, and dedupes on
\code{StripeWebhookEvent.StripeEventId} so a re-delivered event is processed once.

\begin{secbox}
Three independent guarantees protect the money path: \textbf{idempotency keys}
(no double-charge / double-payout on retries), \textbf{manual capture} (funds are
only taken when the service is actually committed), and \textbf{signed webhooks +
event dedupe} (no spoofed or replayed payment events). Card details never touch
our servers --- the browser talks to Stripe directly via the client secret; we
only ever see the intent id.
\end{secbox}

\begin{qabox}
``Where is the money between booking and the session?'' --- It is \emph{authorized
and held} on the student's card via a manual-capture PaymentIntent. We capture it
only when the consultant accepts; if they reject or the slot expires we release the
hold; if there's a qualifying no-show/cancellation the RefundCalculatorService
decides the refund. Net earnings reach the consultant through a Stripe Connect
payout.
\end{qabox}

% ====================================================================
\section{Cross-cutting engineering (quick hits)}
\begin{itemize}
\item \textbf{Audit log} --- \code{AuditService} writes append-only \code{AuditLog}
  rows (before/after JSON, polymorphic target).
\item \textbf{Real-time} --- three SignalR hubs (Chat, Notification, Community).
\item \textbf{Background jobs} --- 13 Hangfire jobs (reminders, scholarship
  auto-close, no-show sweep, payouts, GDPR export/delete, redaction sampling,
  integrity checks), \emph{feature-flagged} via \code{Hangfire:Enabled}.
\item \textbf{Bilingual} EN/AR with RTL throughout.
\item \textbf{Documentation rigor} --- EERD, Relational-Mapping, Class and Component
  reports, each gap-analyzed against the live code and corrected.
\end{itemize}

\vspace{6pt}\hrule height 0.4pt\vspace{4pt}
{\small\itshape\color{spInk} All mechanisms above are implemented in
\code{server/src/...} and verified against the source for this brief.}

\end{document}
